% GENERATED by scripts/build-cv.mjs from data/portfolio-data.json. Edit the data (or the CMS), not this file.
\documentclass[10pt]{article}
\usepackage[a4paper,margin=1.7cm,top=1.4cm,bottom=1.4cm]{geometry}
\input{glyphtounicode}
\pdfgentounicode=1
\usepackage[T1]{fontenc}
\usepackage[utf8]{inputenc}
\usepackage{XCharter}
\usepackage[none]{hyphenat}
\usepackage{ragged2e}
\usepackage{microtype}
\usepackage{enumitem}
\usepackage{titlesec}
\usepackage{xcolor}
\usepackage[hidelinks,pdfusetitle]{hyperref}
\usepackage{xurl}
\urlstyle{same}
\hypersetup{pdftitle={Durjoy Majumdar -- Curriculum Vitae}, pdfauthor={Durjoy Majumdar}, pdfsubject={Academic CV}, pdfkeywords={Cross-chain \& blockchain security, Agentic AI for security, Adversarial uses of LLMs, Reliability of LLMs in security analysis, Network \& protocol security, Trust assumptions in distributed systems, blockchain security, cross-chain bridges, smart contracts, LLM security, agentic AI, cybersecurity}}
\newcommand{\doi}[1]{\href{https://doi.org/#1}{\nolinkurl{doi:#1}}}
\definecolor{accent}{HTML}{1F3A5F}
\definecolor{muted}{HTML}{555555}
\pagestyle{empty}
\setlength{\parindent}{0pt}
\setlength{\parskip}{0pt}
\linespread{1.0}
\RaggedRight
\titleformat{\section}{\large\bfseries\color{accent}}{}{0pt}{}[\vspace{-0.55em}\textcolor{accent}{\rule{\linewidth}{0.5pt}}]
\titlespacing*{\section}{0pt}{0.7em}{0.3em}
\setlist[itemize]{leftmargin=1.1em,itemsep=0.06em,topsep=0.12em,parsep=0pt,label={\textcolor{muted}{\textbullet}}}
\newcommand{\entry}[3]{%
  \textbf{#1}\hfill{\color{muted}#2}\par
  \ifx&#3&\else{\itshape #3}\par\fi\vspace{-0.1em}}
\begin{document}
\begin{center}
  {\LARGE\bfseries Durjoy Majumdar}\\[0.3em]
  M.Tech, Computer Science \& Engineering, IIT Patna | Junior Researcher, Cybersecurity Lab\\[0.25em]
  {\small \href{mailto:durjoy_2511cs25@iitp.ac.in}{durjoy\_2511cs25@iitp.ac.in} \enspace{}|\enspace{} \href{mailto:durjoymajumdar02@gmail.com}{durjoymajumdar02@gmail.com} \enspace{}|\enspace{} +91 78109 76687 \enspace{}|\enspace{} \href{https://orcid.org/0009-0002-3086-9893}{ORCID 0009-0002-3086-9893}\\
  \href{https://asokakrsna.github.io}{asokakrsna.github.io} \enspace{}|\enspace{} \href{https://github.com/AsokaKrsna}{github.com/AsokaKrsna} \enspace{}|\enspace{} \href{https://www.linkedin.com/in/durjoy-majumdar/}{linkedin.com/in/durjoy-majumdar} \enspace{}|\enspace{} \href{https://scholar.google.com/citations?user=c7OEFoIAAAAJ}{Google Scholar}}
\end{center}
\vspace{-0.5em}

\section{Research Interests}
How systems fail when their trust assumptions meet an adversary who does not share them: blockchain and cross-chain bridge security; smart contract and DeFi security; agentic AI and large language models (LLMs) for security analysis and attack detection; adversarial uses of LLMs; network security and protocol-level attacks.

\section{Education}
\entry{Indian Institute of Technology Patna}{2025 -- 2027 (expected)}{M.Tech, Computer Science \& Engineering \hfill CPI 8.73/10 (current)}
\begin{itemize}
  \item Supervisor: Dr. Somnath Tripathy. Thesis area: agentic AI for cross-chain security.
\end{itemize}
\vspace{0.3em}
\entry{KL University (KLEF), Vaddeswaram}{2021 -- 2025}{B.Tech (Honours), Computer Science \& Engineering \hfill CGPA 9.19/10}
\begin{itemize}
  \item Specialisation in Cybersecurity and Blockchain Technology.
\end{itemize}

\section{Publications}
\textit{Preprints}
\begin{itemize}
  \item \textbf{D.~Majumdar}. ``Beyond F1: A Study of LLM Reliability in Smart Contract Vulnerability Detection.'' \emph{Preprint, Research Square}, Aug 2026. \doi{10.21203/rs.3.rs-10783572/v1}
\end{itemize}
\vspace{0.2em}
\textit{Peer-reviewed conference papers}
\begin{itemize}
  \item T. R. Smriti, \textbf{D.~Majumdar}. ``Hybrid Machine-Deep Learning Ensembles for English, Bangla and Banglish Spam Text Detection.'' \emph{IEEE Pune Section International Conference (PuneCon), pp. 1--6}, 2025. \doi{10.1109/PuneCon67554.2025.11377913}
  \item \textbf{D.~Majumdar}, S. Arjun, P. Boyina, S. S. P. Rayidi, Y. R. Sai, S. V. Gangashetty. ``Beyond Text: Nefarious Actors Harnessing LLMs for Strategic Advantage.'' \emph{International Conference on Intelligent Systems for Cybersecurity (ISCS), IEEE}, 2024. \doi{10.1109/ISCS61804.2024.10581181}
  \item S. Arjun, \textbf{D.~Majumdar}, A. Z. Nabin, A. K. Sharma, S. M. Raheman. ``Beyond Copy–Pasting—Contextualizing LLMs for Secure Code Generation.'' \emph{International Conference on Innovations in Cybersecurity and Data Science (ICICDS), Springer}, 2024. \doi{10.1007/978-981-97-5791-6\_15}
  \item \textbf{D.~Majumdar}, R. R. Chintala. ``Shadow ARP-ing for ARP Poisoning Detection.'' \emph{9th International Conference on Advanced Computing and Communication Systems (ICACCS), IEEE}, 2023. \doi{10.1109/ICACCS57279.2023.10112712}
\end{itemize}

\section{Research Experience}
\entry{Junior Researcher, Cybersecurity Lab, IIT Patna}{Nov 2025 -- Present}{}
\begin{itemize}
  \item Research on blockchain and cross-chain bridge security, supervised by Dr. Somnath Tripathy.
  \item Built Crucible, a testbed that reproduces historical cross-chain bridge exploits on forked EVM chains, and a private multi-node Ethereum network for smart contract security experiments (see Research Software).
  \item Ongoing: agentic AI for cross-chain attack detection (M.Tech thesis); with a lab summer intern, making hashed time-lock contracts (HTLCs) resistant to miner bribery.
\end{itemize}
\vspace{0.3em}
\entry{Cybersecurity Researcher (part-time), CRAC Learning}{Apr 2025 -- Nov 2026}{}
\begin{itemize}
  \item Studied how adversaries misuse AI models beyond conventional uses, and how AI can support defensive security.
\end{itemize}
\vspace{0.3em}
\entry{Research Associate (Intern), Speech Research Group, KL University}{Nov 2024 -- May 2025}{}
\begin{itemize}
  \item Supported the lab's data management, systems and security, and a project on voice-authentication attacks (replay and deepfake).
\end{itemize}
\vspace{0.3em}
\entry{Student Researcher, AI4ECN Horizon Europe consortium proposal}{}{}
\begin{itemize}
  \item Represented KL University in a Horizon Europe consortium proposal on AI for the edge-cloud-network continuum, with 17 international partners including UCL and Telefónica.
\end{itemize}

\section{Research Software}
\begin{itemize}
  \item \textbf{Crucible} -- Modular, CLI-first testbed for EVM cross-chain bridge security research: forks chains at historical blocks, replays exploits deterministically, records run artifacts with integrity hashes, and offers a plugin SDK for attack, bridge and LLM-agent plugins. Benchmarks for six incidents, including Wormhole, Nomad and Poly Network. Python, Foundry, Solidity, Docker. \href{https://asokakrsna.github.io/crucible\_docs/}{\nolinkurl{asokakrsna.github.io/crucible\_docs}}
  \item \textbf{Private Ethereum Research Testbed} -- Multi-node proof-of-work Ethereum network on lab machines at IIT Patna, with a faucet and role-based dashboard. Geth, Docker, Flask, Nginx.
  \item \textbf{CryptoSIM} -- Desktop application for teaching cryptography and blockchain through simulation; 310 primitives and an instructor Professor Mode. Rust, Tauri. \href{https://cryptosim-a2e.pages.dev/}{\nolinkurl{cryptosim-a2e.pages.dev}}
\end{itemize}

\section{Projects}
\begin{itemize}
  \item \textbf{0xProblems} -- Atlas of 229 open cybersecurity research problems across 11 areas, for students choosing a research direction. JavaScript. \href{https://0xproblems.vercel.app/}{\nolinkurl{0xproblems.vercel.app}}
  \item \textbf{Logalyzer} -- SIEM for real-time log analysis and automated response, with log collection from isolated endpoints. Smart India Hackathon 2023 finalist. Python, PHP, PowerShell, JavaScript. \href{https://github.com/AsokaKrsna/SIH2023}{\nolinkurl{github.com/AsokaKrsna/SIH2023}}
  \item \textbf{WeCrypt} -- Forensic tool to locate and decrypt WeChat message databases via a downgrade weakness, tested on 50+ databases. KAVACH 2023 finalist. JavaScript, Shell. \href{https://github.com/AsokaKrsna/KAVACH\_2023}{\nolinkurl{github.com/AsokaKrsna/KAVACH\_2023}}
\end{itemize}

\section{Honours and Awards}
\begin{itemize}[leftmargin=4.8em,labelwidth=4.3em,labelsep=0.5em,align=left]
  \item[{\color{muted}2026}] Selected to attend the Google DeepMind APAC Research Symposium (Student Days).
  \item[{\color{muted}2025--27}] ICCR scholarship, Government of India, for M.Tech study at IIT Patna.
  \item[{\color{muted}2025}] Conference grant, Seasides Cybersecurity Conference.
  \item[{\color{muted}2023}] Finalist, Smart India Hackathon and KAVACH national cybersecurity hackathon (top 5 teams nationally per problem statement).
  \item[{\color{muted}2021--25}] Study in India (SII) scholarship, Government of India; 250 awardees worldwide, selected by the PRAGATII exam.
  \item[{\color{muted}2019}] National finalist, Extraordinary Talent Hunt (Science), Bangladesh.
\end{itemize}

\section{Teaching and Mentoring}
\entry{Teaching Assistant, IIT Patna}{Aug 2025 -- Present}{}
\begin{itemize}
  \item Teaching assistant for Natural Language Processing and Advanced Topics in Blockchain.
\end{itemize}
\vspace{0.3em}
\entry{Peer Mentor, KL University}{Apr 2022 -- Jun 2024}{}
\begin{itemize}
  \item Mentored 100+ undergraduates across 7 security courses, including Computer Networks \& Security and Digital Forensics; co-developed 4 lab workbooks used by 2,000+ students with 8 faculty members.
\end{itemize}
\vspace{0.3em}
\entry{Volunteer Teaching and Mentoring}{}{}
\begin{itemize}
  \item Free cybersecurity awareness sessions for underprivileged students in Bengali (2024 – Present); volunteer mentor with Skillmate and Gold Heart Foundation (2025 – Present).
\end{itemize}

\section{Service and Leadership}
\entry{Research Advisor, White Hat Hackers Club, KL University}{May 2024 -- Apr 2025}{}
\begin{itemize}
  \item Initiated 6 student research projects and guided 10 members, leading to 2 published papers; ran biweekly research discussions.
\end{itemize}
\vspace{0.3em}
\entry{Technical Chairperson, White Hat Hackers Club, KL University}{May 2023 -- May 2024}{}
\begin{itemize}
  \item Coordinated 15+ security events with 1,500+ participants; authored CTF challenges for 4 events, many built around covert channels.
\end{itemize}
\vspace{0.3em}
\entry{CTF Player, Team Enigma Hunters}{2022 -- Present}{}
\begin{itemize}
  \item Competitive capture-the-flag player; top 1\% globally on TryHackMe and LetsDefend.
\end{itemize}

\section{Skills}
\begin{itemize}
  \item \textbf{Programming:} Python, Rust, Solidity, C, JavaScript, Bash, PowerShell, LaTeX; Linux, Docker, Git
  \item \textbf{Blockchain:} Ethereum, EVM, smart contracts, cross-chain bridges, Foundry (Forge, Anvil, Cast), Geth, chain forking and transaction replay
  \item \textbf{AI and LLMs:} LLM evaluation and benchmarking, prompt engineering, LLM agents, Jupyter, pandas
  \item \textbf{Security:} Vulnerability analysis, incident response, log and malware analysis, digital forensics, intrusion detection; Wireshark, Burp Suite, Nmap, Metasploit, Splunk, Volatility, Snort, Suricata
  \item \textbf{Certifications:} ISC2 Certified in Cybersecurity (CC); Juniper JNCIA-Junos; Fortinet Certified Associate in Cybersecurity; Google Cybersecurity Certificate; OCI Generative AI Professional; TryHackMe SAL1
\end{itemize}
\end{document}
